\scanpage{119}
de $T$ défini par
\[
\unsure{D|\mathbb D}=V^\perp\tag*{(3)}
\]
et en notant que
\[
D\simeq(E/V)^\vee\simeq\mathcal O_S^\vee\simeq\mathcal O_S\tag*{(4)}
\]
i.e. $D$ est muni canoniquement d'une \deleted{section}\inserted{base} $a$, \unsure{d.à.d.|i.e.} d'une section
\[
\left\{\begin{aligned}
&a\in\Gamma(S,T)\deleted{\unsure{\subset\Gamma(|\subset\Gamma(T)}}\\
&E=\{x\in\Gamma(-,T)\mid\varphi(x,a)=1\}
\end{aligned}\right.\tag*{(4')}
\]
-- du moins quand $q$ est \uline{\unsure{choisie|définie|donnée}} \deleted{\unsure{sans}} (non seulement mod une unité). Dans le cas général, regardant $q$ comme forme à valeurs dans un module inv. $L$ (quotient de $\Gamma^2(T)$, i.e. dual d'un ss-fibré inv. de $\operatorname{Sym}^2(\check T)\ldots$) on trouve $D\otimes(\unsure{T|E}/V)\simeq L$ i.e. [comme $T/V\simeq\mathcal O_S$] on trouve un isom canonique
\[
D\xrightarrow{\sim}L\tag*{(5)}
\]
La condition $H$ non tangent à $Q$ s'exprime alors par la condition que
% archive sheet number: 118
